\section{Rational collective readouts with finite risk certificates}
\label{sec:finiterisk81}

The effective selection of Section~\ref{sec:effectivelearning80}
can be realized using rational arithmetic and finite enumeration.
The readout operators themselves may be chosen Gaussian rational.
Their uniform guarantee is certified by exact inequalities on an
explicit finite net.  The net certifies the statistical risk of a
fixed readout; it does not discretize the unknown measurement in
the learning model.

\subsection{A finite certificate for a collective readout}
Keep the public dictionary
$\mathcal D=(C_1,\ldots,C_L)$ of
Theorem~\ref{thm:interiorcodec79}, constructed at $(d,N,\delta)$.
Put $k=\lceil\sqrt N\rceil$ and $a=\delta/(8k)$.
For $m\ge1$, let $D=d^m$ and retain the subnormalized Choi records
$\Gamma_m(E)$ of \eqref{eq:choirecord80}.
For rational $0<\eta\le1/8$, define
\begin{equation}\label{eq:certificateparameters81}
 r_m=\min\left\{\frac{\delta}{32k},\frac{\eta}{16m}\right\},
 \qquad
 K_m=\left\lceil\frac{4d}{r_m}\right\rceil,
 \qquad
 Q_m=\left\lceil\frac{64L^2D}{\eta}\right\rceil.
\end{equation}
Let $\mathcal G_m$ be the ordered finite set of Gaussian-rational
Hermitian matrices with denominator $K_m$ in
$\mathfrak I_d=[I_d/4,3I_d/4]$.  More explicitly, enumerate diagonal
coordinates in $\{0,1/K_m,\ldots,1\}$ and real and imaginary
coordinates above the diagonal in
$\{-1,-1+1/K_m,\ldots,1\}$, in the public order, and retain exactly
those matrices in $\mathfrak I_d$.
All tests use exact rational positive-semidefinite elimination.
For $G\in\mathcal G_m$, set
\[
 J(G)=\{j:\|G-C_j\|_{\rm op}\le a\}.
\]
These sets are decidable by the two positive-semidefinite tests in
\eqref{eq:interiortest79}, including equality.

A rational readout is a family
$(A_{j,y}:1\le j\le L,\ y\in\{0,1\}^m)$ of Gaussian-rational
Hermitian $D\times D$ matrices satisfying
\[
 A_{j,y}\succeq0,\qquad \sum_{j=1}^L A_{j,y}=I_D
 \quad\hbox{for every }y.
\]
Its certificate consists of these exact legality conditions and the
finite inequalities
\begin{equation}\label{eq:finitecertificate81}
 \sum_{j\in J(G)}\sum_{y\in\{0,1\}^m}
       \operatorname{tr}(A_{j,y}R_y(G))\ge1-\frac\eta4
 \qquad(G\in\mathcal G_m).
\end{equation}
Every term is rational.  A verifier reconstructs the dictionary and
net from the public parameters and checks every displayed inequality.
A proper prefix of either enumeration does not constitute a certificate.

\begin{theorem}[Uniform validity of the finite certificate]
\label{thm:finitecertificate81}
Every readout satisfying \eqref{eq:finitecertificate81} has, for all
real $E\in\mathfrak I_d$,
\begin{equation}\label{eq:certifiedrisk81}
 \mathbb P_E\{d_N(E,C_J)>5\delta/8\}\le5\eta/16.
\end{equation}
Here $J$ is the readout index, and $C_J$ is its unchanged public
rational decoding.  More generally, if a finite operator-norm $r$-net
$\mathcal G\subseteq\mathfrak I_d$ satisfies the same finite inequalities with right side
$1-\alpha$ and good-label radius $a$, then
\begin{equation}\label{eq:certificatetransfer81}
 \sup_{E\in\mathfrak I_d}\mathbb P_E\{\|E-C_J\|_{\rm op}>a+r\}
 \le\alpha+mr.
\end{equation}
The centres may be any fixed finite legal family for this last assertion.
\end{theorem}
\begin{proof}
We first verify that $\mathcal G_m$ is an operator-norm $r_m$-net.
Write $K=K_m$, and for any $E\in\mathfrak I_d$ form
\[
 T=\left(1-\frac{8d}{K}\right)E+\frac{4d}{K}I_d.
\]
Since $8d/K\le2r_m<1$, this is a convex displacement towards
$I_d/2$.  Its spectrum lies in
$[1/4+2d/K,3/4-2d/K]$ and
$\|T-E\|_{\rm op}\le2d/K$.
Round each independent real coordinate to the nearest multiple of
$1/K$, with the prescribed ties, and preserve Hermitian conjugacy.
The row-sum estimate gives $\|G-T\|_{\rm op}\le d/K$.
Consequently $G\in\mathfrak I_d$, its coordinates occur in the
specified finite stream, and
\begin{equation}\label{eq:certificatenet81}
 \|G-E\|_{\rm op}\le3d/K\le3r_m/4\le r_m.
\end{equation}

For one Choi record,
\[
 \|\Gamma_1(E)-\Gamma_1(G)\|_1
   =\frac{2}{d}\|E-G\|_1\le2\|E-G\|_{\rm op}.
\]
Tensor-product telescoping between normalized states therefore yields
\begin{equation}\label{eq:recordlipschitz81}
 \|\Gamma_m(E)-\Gamma_m(G)\|_1
 \le2m\|E-G\|_{\rm op}.
\end{equation}
Apply the same fixed readout under the two target effects.
Their index laws have total variation at most $mr_m\le\eta/16$.
For the particular $G$ in \eqref{eq:certificatenet81}, every
$j\in J(G)$ satisfies
\[
 \|E-C_j\|_{\rm op}\le a+r_m
                  \le\frac{5\delta}{32k}.
\]
Both effects lie in $[I_d/8,7I_d/8]$, so
Lemma~\ref{lem:interiormetric78} gives
$d_N(E,C_j)\le4\sqrt N(a+r_m)\le5\delta/8$.
The event $\{J\in J(G)\}$ has probability at least
$1-\eta/4-\eta/16$ under $E$, by
\eqref{eq:finitecertificate81}.  This proves
\eqref{eq:certifiedrisk81}.  No continuity of the index-selection
boundary is used: one fixed subset $J(G)$ is transferred between the
two experiments.  Replacing the numerical allowances by $a,r,\alpha$
in that same argument proves \eqref{eq:certificatetransfer81}.
\end{proof}

\subsection{A prescribed rational denominator suffices}
\begin{proposition}[Rational approximation of a finite classical readout]
\label{prop:rationalreadout81}
Let $(B_{j,y})$ be any family of $L$-outcome POVMs on dimension $D$,
with finitely many classical conditioning values $y$.  For rational
$0<\eta\le1/8$, put
$s=\eta/16$ and $Q=\lceil64L^2D/\eta\rceil$.
There is a family $(A_{j,y})$ of Gaussian-rational POVMs, with all real and imaginary
coordinates in $Q^{-1}\mathbb Z$, such that, for every input
classical--quantum state, its output-index law differs from that of
$(B_{j,y})$ in total variation by at most $3\eta/32$.
The rational family belongs to an explicitly finite enumeration
determined solely by $(D,L,Q)$ and the conditioning alphabet.
\end{proposition}
\begin{proof}
When $L=1$, take $A_{1,y}=I_D$; there is no approximation error.
Suppose $L\ge2$.  First put
\[
 \overline B_{j,y}=(1-s)B_{j,y}+\frac{s}{L}I_D.
\]
These effects sum to $I_D$ and are all at least $(s/L)I_D$.
For each $y$, round the independent real coordinates of the first
$L-1$ effects to multiples of $1/Q$, preserving Hermitian conjugacy,
and call the rounded matrices $A_{j,y}$.  Set
\[
 A_{L,y}=I_D-\sum_{j<L}A_{j,y}.
\]
Thus normalization holds exactly.  With
$F_{j,y}=A_{j,y}-\overline B_{j,y}$, the row-sum estimate gives
\[
 \|F_{j,y}\|_{\rm op}\le D/Q\quad(j<L),\qquad
 \|F_{L,y}\|_{\rm op}\le(L-1)D/Q.
\]
Since
\[
 (L-1)D/Q\le\frac{\eta(L-1)}{64L^2}
              <\frac{\eta}{16L}=\frac{s}{L},
\]
all $A_{j,y}$ are positive.  Together with their exact sum, this
also implies $A_{j,y}\preceq I_D$.  The rounded coordinates therefore
lie in the finite ranges $[0,1]$ on the diagonal and $[-1,1]$ for
each real or imaginary off-diagonal coordinate.

Mixing with the uniform outcome changes every event probability by
at most $s$.  Rounding changes it by at most
\[
 \sum_j\|F_{j,y}\|_{\rm op}
 \le2(L-1)D/Q\le2LD/Q\le\eta/32.
\]
The same bound holds after averaging over any subnormalized
classical conditioning blocks.  Thus the total variation of the
index laws is at most $s+\eta/32=3\eta/32$.

For clarity, no exact entries of $B$ are required by the algorithm
below.  This rounding argument proves that a suitable rational tuple
exists in a prescribed finite set.  Enumerate the first $L-1$
Hermitian matrices in the indicated coordinate ranges, define the
last by the residual formula, and retain only positive tuples.
The resulting finite stream contains the family just constructed.
\end{proof}

\begin{theorem}[Certified effective learning at the joint rates]
\label{thm:rationallearner81}
Let $d,N\ge1$, and let $0<\delta\le2^{-13}$ and
$0<\eta\le1/8$ be rational.  A finite deterministic computation
from $(d,N,\delta,\eta)$ returns a query number $M$, a
Gaussian-rational collective readout, and the finite risk certificate
\eqref{eq:finitecertificate81}.  In the finite trusted-control model,
the resulting common learner on $\mathfrak I_d$ satisfies
\begin{align}
 \sup_{E\in\mathfrak I_d}
       \mathbb P_E\{d_N(E,C)>\delta\}&\le\eta,
                                         \label{eq:certifiedactual81}\\
 M&\le C_0\frac{N}{\delta^2}
              \left(d^2+d\log\frac1\eta\right),
                                         \label{eq:certifiedcalls81}\\
 B&=\frac{d^2}{2}\log_2N+d^2\log_2(1/\delta)+O(d^2).
                                         \label{eq:certifiedbits81}
\end{align}
All constants are absolute.  The acquisition consists of fresh
one-call probe--reference pairs, followed by collective processing
of completed outputs.  The call and description orders are the
minimax orders of Theorems~\ref{thm:interiorconfidence80} and
\ref{thm:interiorentropy79}.  No classical running-time, workspace,
or elementary known-gate bound is asserted.
\end{theorem}
\begin{proof}
Construct the public dictionary once.  For each
$m=1,2,4,\ldots$, compute \eqref{eq:certificateparameters81} and
exhaust the finite rational readout stream of
Proposition~\ref{prop:rationalreadout81}, with $2^m$ conditioning values
and dimension $D=d^m$.  For every candidate, check exact POVM
legality and every inequality in \eqref{eq:finitecertificate81}.
If one passes, return the first such candidate and set $M=m$.
If none passes, the finite enumeration ends and the algorithm moves
to $2m$.  Thus it does not wait indefinitely at an infeasible query
count.  All arithmetic, tensor products, comparisons, and
positive-semidefinite tests are over the Gaussian rationals.

It remains to show that some stage passes at the claimed call order.
Lemma~\ref{lem:binaryconfidenceupper80}, used at accuracy
$\epsilon=\delta/(128k)$ and failure $\eta/8$, supplies an
independent-Choi estimator with legal output $F$ and
\[
 m_0\le C_1 k^2\delta^{-2}
                  \left(d^2+d\log\frac8\eta\right).
\]
Clip $F$ spectrally to $[I_d/4,3I_d/4]$, obtaining $T$.
On its good event, Weyl's inequality gives
$\|T-E\|_{\rm op}\le2\epsilon$.
The dictionary rounding and predecessor selection give
\[
 \|T-C_j\|_{\rm op}\le\frac{5\delta}{64k},\qquad
 \|E-C_j\|_{\rm op}\le\frac{3\delta}{32k}<a.
\]
These operations define a Borel finite partition of the original
tomography outcomes.  Hence there exists an $L$-outcome collective
POVM, dephased in the classical acquisition strings, with failure
at most $\eta/8$ for the operator-radius-$a$ event uniformly on
$\mathfrak I_d$.

For every $m\ge m_0$, this readout is available on $m$ Choi records
by ignoring the last $m-m_0$ records.  Equivalently, its effects are
tensored with the identity on the extra reference factors and do
not depend on the extra outcome bits.  Apply
Proposition~\ref{prop:rationalreadout81} at this $m$.
At every $G\in\mathcal G_m$, the resulting rational candidate has
failure at most
\[
 \frac\eta8+\frac\eta{16}+\frac\eta{32}
   =\frac{7\eta}{32}<\frac\eta4.
\]
It therefore passes every finite certificate inequality.
The first dyadic integer at least $m_0$ is at most $2m_0$; thus the
outer search terminates with $M\le2m_0$.  This proof uses existence
of the tomography readout only.  Its entries, continuous outcome
integrals, and exact optimization outputs are never inputs to the
rational search.  Since $k^2\le4N$ and
$\log(8/\eta)=O(\log(1/\eta))$, the query bound follows.

By Theorem~\ref{thm:finitecertificate81}, the selected exact readout
has original-loss failure at most $5\eta/16$, with successful loss
at most $5\delta/8$.  Each rational POVM has an effectively
computable algebraic Stinespring isometry, obtained from the
positive square roots of its effects as in
Section~\ref{sec:effectivelearning80}.  Use the normalized rational
approximations of Section~\ref{sec:finitecontrol78} to specify its
trusted implementation.  There are $M$ fresh-pair preparation
events and one complete classically controlled readout event.
Give each event total unhalved error allowance $\eta/(M+1)$,
split equally between finite target approximation and actual
trusted realization.  This final event includes the complete
readout, its classical conditioning, and its final dephasing.
If it is implemented as several separately charged instructions,
subdivide its allowance among those instructions.

The adaptive control-error lemma bounds the total variation of the
ideal and actual index laws by $\eta/2$.  The exact public decoder
is unchanged.  Consequently
\[
 \mathbb P_E\{d_N(E,C)>\delta\}
 \le\frac{5\eta}{16}+\frac\eta2
 =\frac{13\eta}{16}<\eta
 \qquad(E\in\mathfrak I_d).
\]
Approximate pair preparations remain tensor-separated from every
stored reference, so the fresh one-call restriction is preserved.
The dictionary index supplies the description order in
\eqref{eq:certifiedbits81}; its lower order follows from the
interior covering converse.  The unrestricted adaptive learning
converse applies to this learner as well.
\end{proof}

The certificate consists entirely of finite exact rational checks.
Its continuum implication is Theorem~\ref{thm:finitecertificate81},
whose analytic input is the tensor-product Lipschitz estimate and
the dimension-uniform interior metric comparison.  Exhaustive
rational search replaces real quantifier elimination in this
construction.  The output arrays have $2^M Ld^{2M}$ independent real matrix
coordinates before normalization is used; their
construction and verification are separate resources from the
transmitted dictionary index.  At stage $m$, at most
$(2K_m+1)^{d^2}$ coordinate tuples are inspected for the parameter
net, and at most $(2Q_m+1)^{(L-1)2^mD^2}$ tuples for the readout.
At a fixed stage $m$, the readout array can be stored using $(L-1)2^mD^2\lceil\log_2(2Q_m+1)\rceil$
bits in its fixed coordinate order, together with the public parameters;
the residual effects are reconstructed exactly.  The stage $m$ is
included in a finite header or recovered by rerunning the prescribed
deterministic search.  These explicit finite
bounds are not polynomial-time or payload bounds.  A software check of a finite
instance establishes that instance's certificate, while the stated
lemmas establish its uniform interpretation and the termination
bound for the complete construction.
